% paper74-copilot-construction-participant.tex
%   CopilotConstructionParticipant — LLM-guided Proof Construction
%   Paper 74 of the Judgment Geometry series.
% Compile: pdflatex paper74-copilot-construction-participant
\documentclass[11pt]{article}
\usepackage{lmodern}
% jugeo-common.tex — shared preamble for all Judgment Geometry papers
% Include via: % jugeo-common.tex — shared preamble for all Judgment Geometry papers
% Include via: % jugeo-common.tex — shared preamble for all Judgment Geometry papers
% Include via: \input{jugeo-common}

% ── Packages ────────────────────────────────────────────────────────────
\usepackage{amsmath,amssymb,amsthm,mathtools}
\usepackage{tikz-cd}
\usepackage{listings}
\usepackage{xcolor}
\usepackage{xspace}
\usepackage{booktabs}
\usepackage{multirow}
\usepackage{enumitem}
\usepackage[expansion=false]{microtype}
\usepackage{hyperref}
\usepackage{cleveref}
% \usepackage{thmtools}  % disabled: not available in this TeX installation
\usepackage{float}
\usepackage{caption}
\usepackage{subcaption}
\usepackage{tabularx}
\usepackage{stmaryrd}       % \llbracket, \rrbracket
\usepackage{bbm}            % \mathbbm{1}

% ── Page geometry ───────────────────────────────────────────────────────
\usepackage[margin=1in]{geometry}

% ── Theorem environments ────────────────────────────────────────────────
\theoremstyle{plain}
\newtheorem{theorem}{Theorem}[section]
\newtheorem{lemma}[theorem]{Lemma}
\newtheorem{proposition}[theorem]{Proposition}
\newtheorem{corollary}[theorem]{Corollary}

\theoremstyle{definition}
\newtheorem{definition}[theorem]{Definition}
\newtheorem{example}[theorem]{Example}
\newtheorem{construction}[theorem]{Construction}

\theoremstyle{remark}
\newtheorem{remark}[theorem]{Remark}
\newtheorem{notation}[theorem]{Notation}

% ── Python listings style ──────────────────────────────────────────────
\definecolor{jugeoBlue}{HTML}{2563EB}
\definecolor{jugeoGreen}{HTML}{059669}
\definecolor{jugeoRed}{HTML}{DC2626}
\definecolor{jugeoPurple}{HTML}{7C3AED}
\definecolor{jugeoGray}{HTML}{6B7280}
\definecolor{jugeoBg}{HTML}{F8FAFC}

\lstdefinestyle{jugeo-python}{
  language=Python,
  basicstyle=\ttfamily\small,
  keywordstyle=\color{jugeoBlue}\bfseries,
  stringstyle=\color{jugeoGreen},
  commentstyle=\color{jugeoGray}\itshape,
  numberstyle=\tiny\color{jugeoGray},
  backgroundcolor=\color{jugeoBg},
  numbers=left,
  numbersep=6pt,
  frame=single,
  framerule=0.4pt,
  rulecolor=\color{jugeoGray!40},
  breaklines=true,
  breakatwhitespace=true,
  showstringspaces=false,
  tabsize=4,
  xleftmargin=1.5em,
  framexleftmargin=1.5em,
  morekeywords={dataclass,frozen,field,Enum,IntEnum,auto,Any,
                Mapping,Sequence,Optional,match,case},
  literate={→}{{$\to$}}1 {←}{{$\leftarrow$}}1
           {≤}{{$\leq$}}1 {≥}{{$\geq$}}1 {≠}{{$\neq$}}1
           {∈}{{$\in$}}1 {∀}{{$\forall$}}1 {∃}{{$\exists$}}1
           {⊕}{{$\oplus$}}1 {⊖}{{$\ominus$}}1,
  escapeinside={(*@}{@*)},
}
\lstset{style=jugeo-python}

\lstdefinestyle{jugeo-output}{
  basicstyle=\ttfamily\small,
  backgroundcolor=\color{jugeoBg},
  frame=single,
  framerule=0.4pt,
  rulecolor=\color{jugeoGray!40},
  breaklines=true,
  showstringspaces=false,
  xleftmargin=1.5em,
  framexleftmargin=0.5em,
  numbers=none,
}

% ── JuGeo-specific macros ──────────────────────────────────────────────

% Judgment 8-tuple
\newcommand{\J}{\mathcal{J}}
\newcommand{\Jud}[8]{(#1,\,#2,\,#3,\,#4,\,#5,\,#6,\,#7,\,#8)}
\newcommand{\JudShort}[1]{\J_{#1}}

% Components
\newcommand{\coord}{\mathsf{c}}            % coordinate
\newcommand{\prop}{\varphi}                 % proposition
\newcommand{\carrier}{\mathcal{A}}          % carrier type
\newcommand{\evid}{\mathcal{E}}             % evidence bundle
\newcommand{\oblig}{\mathcal{O}}            % obligations
\newcommand{\obstr}{\mathcal{B}}            % obstructions
\newcommand{\trust}{\mathcal{T}}            % trust annotation
\newcommand{\prov}{\Pi}                     % provenance

% Trust algebra
\newcommand{\Talg}{\mathfrak{T}}            % trust ordered algebra
\newcommand{\Eadm}{\mathcal{E}_{\mathrm{adm}}}
\newcommand{\tleq}{\preceq}                 % trust ordering
\newcommand{\tjoin}{\oplus}                 % conservative join
\newcommand{\tatten}{\ominus}               % attenuation
\newcommand{\tpromote}{\uparrow_\pi}        % promotion
\newcommand{\tdemote}{\downarrow_\chi}       % demotion

% Trust tiers
\newcommand{\Tcontra}{\ensuremath{\mathsf{CONTRADICTED}}}
\newcommand{\Tunver}{\ensuremath{\mathsf{UNVERIFIED}}}
\newcommand{\Tcopilot}{\ensuremath{\mathsf{COPILOT}}}
\newcommand{\Toracle}{\ensuremath{\mathsf{ORACLE}}}
\newcommand{\Truntime}{\ensuremath{\mathsf{RUNTIME}}}
\newcommand{\Tsolver}{\ensuremath{\mathsf{SOLVER}}}
\newcommand{\Tproof}{\ensuremath{\mathsf{PROOF}}}

% Site and geometry
\newcommand{\Site}{\mathcal{S}}             % semantic site
\newcommand{\Cov}{\mathrm{Cov}}            % covering family
\newcommand{\Ob}{\mathrm{Ob}}              % objects
\newcommand{\Mor}{\mathrm{Mor}}            % morphisms
\newcommand{\res}{\mathrm{res}}            % restriction
\newcommand{\Cech}{\check{C}}              % Čech complex
\newcommand{\Hcoh}[1]{H^{#1}}             % cohomology group

% Descent
\newcommand{\descent}{\mathrm{desc}}
\newcommand{\glue}{\mathrm{glue}}
\newcommand{\overlap}[2]{#1 \cap #2}

% Semantic moves
\newcommand{\VERIFY}{\textsc{Verify}}
\newcommand{\CONSTRUCT}{\textsc{Construct}}
\newcommand{\NORMALIZE}{\textsc{Normalize}}
\newcommand{\GLUE}{\textsc{Glue}}
\newcommand{\RESTRICT}{\textsc{Restrict}}
\newcommand{\TRANSPORT}{\textsc{Transport}}
\newcommand{\REPAIR}{\textsc{Repair}}
\newcommand{\PROMOTE}{\textsc{Promote}}
\newcommand{\CHALLENGE}{\textsc{Challenge}}
\newcommand{\ESCALATE}{\textsc{Escalate}}

% Fragments
\newcommand{\QFLIA}{\texttt{QF\_LIA}}
\newcommand{\QFLRA}{\texttt{QF\_LRA}}
\newcommand{\QFBV}{\texttt{QF\_BV}}
\newcommand{\QFUF}{\texttt{QF\_UF}}

% Misc
\newcommand{\Python}{\textsc{Python}}
\newcommand{\JuGeo}{\textsc{JuGeo}}
\newcommand{\LEAN}{\textsc{Lean}\,4}
\newcommand{\Fstar}{F$^{\star}$}
\newcommand{\Coq}{\textsc{Coq}}
\newcommand{\Dafny}{\textsc{Dafny}}

% Common shortcuts
\newcommand{\ie}{i.e.\@\xspace}
\newcommand{\eg}{e.g.\@\xspace}
\newcommand{\cf}{cf.\@\xspace}
\newcommand{\wrt}{w.r.t.\@\xspace}

% Evaluation shortcuts
\newcommand{\numcases}{300}
\newcommand{\numspec}{100}
\newcommand{\numeq}{100}
\newcommand{\numbug}{100}
\newcommand{\medtime}{6.5\,ms}
\newcommand{\totaltime}{1.949\,s}


% ── Packages ────────────────────────────────────────────────────────────
\usepackage{amsmath,amssymb,amsthm,mathtools}
\usepackage{tikz-cd}
\usepackage{listings}
\usepackage{xcolor}
\usepackage{xspace}
\usepackage{booktabs}
\usepackage{multirow}
\usepackage{enumitem}
\usepackage[expansion=false]{microtype}
\usepackage{hyperref}
\usepackage{cleveref}
% \usepackage{thmtools}  % disabled: not available in this TeX installation
\usepackage{float}
\usepackage{caption}
\usepackage{subcaption}
\usepackage{tabularx}
\usepackage{stmaryrd}       % \llbracket, \rrbracket
\usepackage{bbm}            % \mathbbm{1}

% ── Page geometry ───────────────────────────────────────────────────────
\usepackage[margin=1in]{geometry}

% ── Theorem environments ────────────────────────────────────────────────
\theoremstyle{plain}
\newtheorem{theorem}{Theorem}[section]
\newtheorem{lemma}[theorem]{Lemma}
\newtheorem{proposition}[theorem]{Proposition}
\newtheorem{corollary}[theorem]{Corollary}

\theoremstyle{definition}
\newtheorem{definition}[theorem]{Definition}
\newtheorem{example}[theorem]{Example}
\newtheorem{construction}[theorem]{Construction}

\theoremstyle{remark}
\newtheorem{remark}[theorem]{Remark}
\newtheorem{notation}[theorem]{Notation}

% ── Python listings style ──────────────────────────────────────────────
\definecolor{jugeoBlue}{HTML}{2563EB}
\definecolor{jugeoGreen}{HTML}{059669}
\definecolor{jugeoRed}{HTML}{DC2626}
\definecolor{jugeoPurple}{HTML}{7C3AED}
\definecolor{jugeoGray}{HTML}{6B7280}
\definecolor{jugeoBg}{HTML}{F8FAFC}

\lstdefinestyle{jugeo-python}{
  language=Python,
  basicstyle=\ttfamily\small,
  keywordstyle=\color{jugeoBlue}\bfseries,
  stringstyle=\color{jugeoGreen},
  commentstyle=\color{jugeoGray}\itshape,
  numberstyle=\tiny\color{jugeoGray},
  backgroundcolor=\color{jugeoBg},
  numbers=left,
  numbersep=6pt,
  frame=single,
  framerule=0.4pt,
  rulecolor=\color{jugeoGray!40},
  breaklines=true,
  breakatwhitespace=true,
  showstringspaces=false,
  tabsize=4,
  xleftmargin=1.5em,
  framexleftmargin=1.5em,
  morekeywords={dataclass,frozen,field,Enum,IntEnum,auto,Any,
                Mapping,Sequence,Optional,match,case},
  literate={→}{{$\to$}}1 {←}{{$\leftarrow$}}1
           {≤}{{$\leq$}}1 {≥}{{$\geq$}}1 {≠}{{$\neq$}}1
           {∈}{{$\in$}}1 {∀}{{$\forall$}}1 {∃}{{$\exists$}}1
           {⊕}{{$\oplus$}}1 {⊖}{{$\ominus$}}1,
  escapeinside={(*@}{@*)},
}
\lstset{style=jugeo-python}

\lstdefinestyle{jugeo-output}{
  basicstyle=\ttfamily\small,
  backgroundcolor=\color{jugeoBg},
  frame=single,
  framerule=0.4pt,
  rulecolor=\color{jugeoGray!40},
  breaklines=true,
  showstringspaces=false,
  xleftmargin=1.5em,
  framexleftmargin=0.5em,
  numbers=none,
}

% ── JuGeo-specific macros ──────────────────────────────────────────────

% Judgment 8-tuple
\newcommand{\J}{\mathcal{J}}
\newcommand{\Jud}[8]{(#1,\,#2,\,#3,\,#4,\,#5,\,#6,\,#7,\,#8)}
\newcommand{\JudShort}[1]{\J_{#1}}

% Components
\newcommand{\coord}{\mathsf{c}}            % coordinate
\newcommand{\prop}{\varphi}                 % proposition
\newcommand{\carrier}{\mathcal{A}}          % carrier type
\newcommand{\evid}{\mathcal{E}}             % evidence bundle
\newcommand{\oblig}{\mathcal{O}}            % obligations
\newcommand{\obstr}{\mathcal{B}}            % obstructions
\newcommand{\trust}{\mathcal{T}}            % trust annotation
\newcommand{\prov}{\Pi}                     % provenance

% Trust algebra
\newcommand{\Talg}{\mathfrak{T}}            % trust ordered algebra
\newcommand{\Eadm}{\mathcal{E}_{\mathrm{adm}}}
\newcommand{\tleq}{\preceq}                 % trust ordering
\newcommand{\tjoin}{\oplus}                 % conservative join
\newcommand{\tatten}{\ominus}               % attenuation
\newcommand{\tpromote}{\uparrow_\pi}        % promotion
\newcommand{\tdemote}{\downarrow_\chi}       % demotion

% Trust tiers
\newcommand{\Tcontra}{\ensuremath{\mathsf{CONTRADICTED}}}
\newcommand{\Tunver}{\ensuremath{\mathsf{UNVERIFIED}}}
\newcommand{\Tcopilot}{\ensuremath{\mathsf{COPILOT}}}
\newcommand{\Toracle}{\ensuremath{\mathsf{ORACLE}}}
\newcommand{\Truntime}{\ensuremath{\mathsf{RUNTIME}}}
\newcommand{\Tsolver}{\ensuremath{\mathsf{SOLVER}}}
\newcommand{\Tproof}{\ensuremath{\mathsf{PROOF}}}

% Site and geometry
\newcommand{\Site}{\mathcal{S}}             % semantic site
\newcommand{\Cov}{\mathrm{Cov}}            % covering family
\newcommand{\Ob}{\mathrm{Ob}}              % objects
\newcommand{\Mor}{\mathrm{Mor}}            % morphisms
\newcommand{\res}{\mathrm{res}}            % restriction
\newcommand{\Cech}{\check{C}}              % Čech complex
\newcommand{\Hcoh}[1]{H^{#1}}             % cohomology group

% Descent
\newcommand{\descent}{\mathrm{desc}}
\newcommand{\glue}{\mathrm{glue}}
\newcommand{\overlap}[2]{#1 \cap #2}

% Semantic moves
\newcommand{\VERIFY}{\textsc{Verify}}
\newcommand{\CONSTRUCT}{\textsc{Construct}}
\newcommand{\NORMALIZE}{\textsc{Normalize}}
\newcommand{\GLUE}{\textsc{Glue}}
\newcommand{\RESTRICT}{\textsc{Restrict}}
\newcommand{\TRANSPORT}{\textsc{Transport}}
\newcommand{\REPAIR}{\textsc{Repair}}
\newcommand{\PROMOTE}{\textsc{Promote}}
\newcommand{\CHALLENGE}{\textsc{Challenge}}
\newcommand{\ESCALATE}{\textsc{Escalate}}

% Fragments
\newcommand{\QFLIA}{\texttt{QF\_LIA}}
\newcommand{\QFLRA}{\texttt{QF\_LRA}}
\newcommand{\QFBV}{\texttt{QF\_BV}}
\newcommand{\QFUF}{\texttt{QF\_UF}}

% Misc
\newcommand{\Python}{\textsc{Python}}
\newcommand{\JuGeo}{\textsc{JuGeo}}
\newcommand{\LEAN}{\textsc{Lean}\,4}
\newcommand{\Fstar}{F$^{\star}$}
\newcommand{\Coq}{\textsc{Coq}}
\newcommand{\Dafny}{\textsc{Dafny}}

% Common shortcuts
\newcommand{\ie}{i.e.\@\xspace}
\newcommand{\eg}{e.g.\@\xspace}
\newcommand{\cf}{cf.\@\xspace}
\newcommand{\wrt}{w.r.t.\@\xspace}

% Evaluation shortcuts
\newcommand{\numcases}{300}
\newcommand{\numspec}{100}
\newcommand{\numeq}{100}
\newcommand{\numbug}{100}
\newcommand{\medtime}{6.5\,ms}
\newcommand{\totaltime}{1.949\,s}


% ── Packages ────────────────────────────────────────────────────────────
\usepackage{amsmath,amssymb,amsthm,mathtools}
\usepackage{tikz-cd}
\usepackage{listings}
\usepackage{xcolor}
\usepackage{xspace}
\usepackage{booktabs}
\usepackage{multirow}
\usepackage{enumitem}
\usepackage[expansion=false]{microtype}
\usepackage{hyperref}
\usepackage{cleveref}
% \usepackage{thmtools}  % disabled: not available in this TeX installation
\usepackage{float}
\usepackage{caption}
\usepackage{subcaption}
\usepackage{tabularx}
\usepackage{stmaryrd}       % \llbracket, \rrbracket
\usepackage{bbm}            % \mathbbm{1}

% ── Page geometry ───────────────────────────────────────────────────────
\usepackage[margin=1in]{geometry}

% ── Theorem environments ────────────────────────────────────────────────
\theoremstyle{plain}
\newtheorem{theorem}{Theorem}[section]
\newtheorem{lemma}[theorem]{Lemma}
\newtheorem{proposition}[theorem]{Proposition}
\newtheorem{corollary}[theorem]{Corollary}

\theoremstyle{definition}
\newtheorem{definition}[theorem]{Definition}
\newtheorem{example}[theorem]{Example}
\newtheorem{construction}[theorem]{Construction}

\theoremstyle{remark}
\newtheorem{remark}[theorem]{Remark}
\newtheorem{notation}[theorem]{Notation}

% ── Python listings style ──────────────────────────────────────────────
\definecolor{jugeoBlue}{HTML}{2563EB}
\definecolor{jugeoGreen}{HTML}{059669}
\definecolor{jugeoRed}{HTML}{DC2626}
\definecolor{jugeoPurple}{HTML}{7C3AED}
\definecolor{jugeoGray}{HTML}{6B7280}
\definecolor{jugeoBg}{HTML}{F8FAFC}

\lstdefinestyle{jugeo-python}{
  language=Python,
  basicstyle=\ttfamily\small,
  keywordstyle=\color{jugeoBlue}\bfseries,
  stringstyle=\color{jugeoGreen},
  commentstyle=\color{jugeoGray}\itshape,
  numberstyle=\tiny\color{jugeoGray},
  backgroundcolor=\color{jugeoBg},
  numbers=left,
  numbersep=6pt,
  frame=single,
  framerule=0.4pt,
  rulecolor=\color{jugeoGray!40},
  breaklines=true,
  breakatwhitespace=true,
  showstringspaces=false,
  tabsize=4,
  xleftmargin=1.5em,
  framexleftmargin=1.5em,
  morekeywords={dataclass,frozen,field,Enum,IntEnum,auto,Any,
                Mapping,Sequence,Optional,match,case},
  literate={→}{{$\to$}}1 {←}{{$\leftarrow$}}1
           {≤}{{$\leq$}}1 {≥}{{$\geq$}}1 {≠}{{$\neq$}}1
           {∈}{{$\in$}}1 {∀}{{$\forall$}}1 {∃}{{$\exists$}}1
           {⊕}{{$\oplus$}}1 {⊖}{{$\ominus$}}1,
  escapeinside={(*@}{@*)},
}
\lstset{style=jugeo-python}

\lstdefinestyle{jugeo-output}{
  basicstyle=\ttfamily\small,
  backgroundcolor=\color{jugeoBg},
  frame=single,
  framerule=0.4pt,
  rulecolor=\color{jugeoGray!40},
  breaklines=true,
  showstringspaces=false,
  xleftmargin=1.5em,
  framexleftmargin=0.5em,
  numbers=none,
}

% ── JuGeo-specific macros ──────────────────────────────────────────────

% Judgment 8-tuple
\newcommand{\J}{\mathcal{J}}
\newcommand{\Jud}[8]{(#1,\,#2,\,#3,\,#4,\,#5,\,#6,\,#7,\,#8)}
\newcommand{\JudShort}[1]{\J_{#1}}

% Components
\newcommand{\coord}{\mathsf{c}}            % coordinate
\newcommand{\prop}{\varphi}                 % proposition
\newcommand{\carrier}{\mathcal{A}}          % carrier type
\newcommand{\evid}{\mathcal{E}}             % evidence bundle
\newcommand{\oblig}{\mathcal{O}}            % obligations
\newcommand{\obstr}{\mathcal{B}}            % obstructions
\newcommand{\trust}{\mathcal{T}}            % trust annotation
\newcommand{\prov}{\Pi}                     % provenance

% Trust algebra
\newcommand{\Talg}{\mathfrak{T}}            % trust ordered algebra
\newcommand{\Eadm}{\mathcal{E}_{\mathrm{adm}}}
\newcommand{\tleq}{\preceq}                 % trust ordering
\newcommand{\tjoin}{\oplus}                 % conservative join
\newcommand{\tatten}{\ominus}               % attenuation
\newcommand{\tpromote}{\uparrow_\pi}        % promotion
\newcommand{\tdemote}{\downarrow_\chi}       % demotion

% Trust tiers
\newcommand{\Tcontra}{\ensuremath{\mathsf{CONTRADICTED}}}
\newcommand{\Tunver}{\ensuremath{\mathsf{UNVERIFIED}}}
\newcommand{\Tcopilot}{\ensuremath{\mathsf{COPILOT}}}
\newcommand{\Toracle}{\ensuremath{\mathsf{ORACLE}}}
\newcommand{\Truntime}{\ensuremath{\mathsf{RUNTIME}}}
\newcommand{\Tsolver}{\ensuremath{\mathsf{SOLVER}}}
\newcommand{\Tproof}{\ensuremath{\mathsf{PROOF}}}

% Site and geometry
\newcommand{\Site}{\mathcal{S}}             % semantic site
\newcommand{\Cov}{\mathrm{Cov}}            % covering family
\newcommand{\Ob}{\mathrm{Ob}}              % objects
\newcommand{\Mor}{\mathrm{Mor}}            % morphisms
\newcommand{\res}{\mathrm{res}}            % restriction
\newcommand{\Cech}{\check{C}}              % Čech complex
\newcommand{\Hcoh}[1]{H^{#1}}             % cohomology group

% Descent
\newcommand{\descent}{\mathrm{desc}}
\newcommand{\glue}{\mathrm{glue}}
\newcommand{\overlap}[2]{#1 \cap #2}

% Semantic moves
\newcommand{\VERIFY}{\textsc{Verify}}
\newcommand{\CONSTRUCT}{\textsc{Construct}}
\newcommand{\NORMALIZE}{\textsc{Normalize}}
\newcommand{\GLUE}{\textsc{Glue}}
\newcommand{\RESTRICT}{\textsc{Restrict}}
\newcommand{\TRANSPORT}{\textsc{Transport}}
\newcommand{\REPAIR}{\textsc{Repair}}
\newcommand{\PROMOTE}{\textsc{Promote}}
\newcommand{\CHALLENGE}{\textsc{Challenge}}
\newcommand{\ESCALATE}{\textsc{Escalate}}

% Fragments
\newcommand{\QFLIA}{\texttt{QF\_LIA}}
\newcommand{\QFLRA}{\texttt{QF\_LRA}}
\newcommand{\QFBV}{\texttt{QF\_BV}}
\newcommand{\QFUF}{\texttt{QF\_UF}}

% Misc
\newcommand{\Python}{\textsc{Python}}
\newcommand{\JuGeo}{\textsc{JuGeo}}
\newcommand{\LEAN}{\textsc{Lean}\,4}
\newcommand{\Fstar}{F$^{\star}$}
\newcommand{\Coq}{\textsc{Coq}}
\newcommand{\Dafny}{\textsc{Dafny}}

% Common shortcuts
\newcommand{\ie}{i.e.\@\xspace}
\newcommand{\eg}{e.g.\@\xspace}
\newcommand{\cf}{cf.\@\xspace}
\newcommand{\wrt}{w.r.t.\@\xspace}

% Evaluation shortcuts
\newcommand{\numcases}{300}
\newcommand{\numspec}{100}
\newcommand{\numeq}{100}
\newcommand{\numbug}{100}
\newcommand{\medtime}{6.5\,ms}
\newcommand{\totaltime}{1.949\,s}


% ——— Paper-specific macros (letters only) ———
\newcommand{\CCP}{\mathsf{CopilotConstructionParticipant}}
\newcommand{\ConStep}{\mathrm{constructionStep}}
\newcommand{\ConPlan}{\mathrm{constructionPlan}}
\newcommand{\ConWitness}{\mathrm{witness}}
\newcommand{\ConVerify}{\mathrm{verify}}
\newcommand{\ConTrust}{\mathrm{trustBound}}
\newcommand{\ConSearch}{\mathrm{searchSpace}}
\newcommand{\CopChan}{\mathsf{CopilotChannel}}
\newcommand{\CopBridge}{\mathsf{CopilotAPIBridge}}
\newcommand{\ChanPool}{\mathsf{ChannelPool}}
\newcommand{\ChanMon}{\mathsf{ChannelMonitor}}
\newcommand{\ChanFed}{\mathsf{ChannelFederation}}
\newcommand{\ChanRouter}{\mathsf{ChannelRouter}}
\newcommand{\EvReq}{\mathsf{EvidenceRequest}}
\newcommand{\EvResp}{\mathsf{EvidenceResponse}}
\newcommand{\EvRec}{\mathsf{EvidenceRecord}}
\newcommand{\SuppRoute}{\mathsf{SupportRoute}}

\title{\textbf{CopilotConstructionParticipant:\\
LLM-Guided Proof Construction in Judgment Geometry}}
\author{Halley Young}
\date{Paper~74 --- Judgment Geometry Series}

\begin{document}
\maketitle

% ============================================================
%  Abstract
% ============================================================
\begin{abstract}
Proof construction---the act of building a verified derivation
from axioms and previously established lemmas---is among the
most resource-intensive operations in the \JuGeo{} verification
pipeline.  This paper introduces the $\CCP$, a Copilot-guided
agent that participates in proof construction by suggesting
construction steps, proposing witnesses, and ranking alternative
construction plans.  The participant operates within the
$\Tcopilot$ trust ceiling: its suggestions are proposals
that must be validated by higher-trust channels (solver, runtime)
before installation.  We formalise construction as a category
of partial proofs, model the participant as a trust-bounded
functor, and prove four theorems: construction soundness,
trust-ceiling preservation, search-space reduction, and
witness completeness.  Four \Python{} listings illustrate
the integration with $\CopChan$, $\CopBridge$, $\ChanFed$,
and the $\mathsf{JuGeoAPI}$.  An appendix collects ten
auxiliary lemmas.
\end{abstract}

% ============================================================
%  Section 1 — Introduction
% ============================================================
\section{Introduction}\label{sec:intro}

In the \JuGeo{} framework, the \texttt{CONSTRUCT} operation
builds a proof object for a given proposition at a given
coordinate.  Unlike \texttt{VERIFY}, which merely checks
an existing proof, construction is creative: the system must
find witnesses, choose induction variables, decide case
splits, and assemble the pieces into a coherent derivation.

Human mathematicians excel at this task because they bring
intuition, analogy, and domain knowledge.  Large language
models (LLMs) share some of these capabilities, at least
superficially.  The $\CCP$ harnesses the LLM's suggestive
power while keeping it firmly within the trust lattice.

\paragraph{Contributions.}
\begin{enumerate}
\item A categorical formalisation of proof construction as
      a functor from partial proofs to completed proofs
      (\S\ref{sec:construction}).
\item The design of the $\CCP$ with three phases: plan
      generation, step suggestion, and witness proposal
      (\S\ref{sec:participant}).
\item Four theorems establishing soundness, trust-ceiling
      preservation, search-space reduction, and witness
      completeness (\S\ref{sec:formal}).
\item Four \Python{} listings showing integration with the
      channel and API layers (\S\ref{sec:code}).
\end{enumerate}

\paragraph{Paper outline.}
Section~\ref{sec:background} gives background.
Section~\ref{sec:construction} formalises proof construction.
Section~\ref{sec:participant} presents the participant design.
Section~\ref{sec:code} shows the implementation.
Section~\ref{sec:formal} states and proves the main theorems.
Section~\ref{sec:discussion} discusses limitations.
Section~\ref{sec:related} surveys related work.
Section~\ref{sec:conclusion} concludes.
Appendix~\ref{app:proofs} collects auxiliary proofs.

% ============================================================
%  Section 2 — Background
% ============================================================
\section{Background}\label{sec:background}

\subsection{Trust lattice}

The trust ordering is:
\[
  \Tcontra \;\tleq\; \Tunver \;\tleq\; \Tcopilot
  \;\tleq\; \Toracle \;\tleq\; \Truntime
  \;\tleq\; \Tsolver \;\tleq\; \Tproof.
\]
The construction participant operates at $\Tcopilot$.
Its suggestions are proposals that the kernel validates
by dispatching to the solver or runtime channel.

\subsection{Evidence channels}

The $\CopChan$ queries an LLM and produces evidence at
trust $\Tcopilot$.  The $\mathsf{SolverChannel}$ checks
validity at $\Tsolver$.  The $\mathsf{RuntimeChannel}$
captures witnesses at $\Truntime$.  The $\ChanFed$
collects evidence from all channels and merges it under
the no-escalation invariant.

\subsection{The CONSTRUCT operation}

The $\mathsf{JuGeoAPI}.\texttt{construct}(c, \varphi)$
operation takes a coordinate~$c$ and a proposition~$\varphi$
and returns either a proof object or a list of residual
sub-goals.  The operation is mediated by the channel layer:
the Copilot bridge suggests a construction plan, and the
solver channel validates each step.

\begin{definition}[Proof obligation]\label{def:proofobligation}
A \emph{proof obligation} is a pair $(c, \varphi)$ where
$c$ is a coordinate in the proof tree and $\varphi$ is
a proposition to be proved.
\end{definition}

\begin{definition}[Partial proof]\label{def:partialproof}
A \emph{partial proof} is a tree whose leaves are either
completed sub-proofs or open goals (residual obligations).
A partial proof with no open goals is a \emph{complete proof}.
\end{definition}

% ============================================================
%  Section 3 — Proof Construction
% ============================================================
\section{Proof Construction}\label{sec:construction}

\subsection{The category of partial proofs}

\begin{definition}[Partial proof category]\label{def:ppcat}
Let $\mathbf{PP}$ be the category whose objects are partial
proofs and whose morphisms are refinement steps: a morphism
$P \to P'$ closes one or more open goals in~$P$ to produce~$P'$.
\end{definition}

\begin{definition}[Refinement step]\label{def:refstep}
A \emph{refinement step} is a function
$r \colon \mathsf{Goal} \to \mathcal{P}(\mathsf{Goal})$
that replaces one open goal with zero or more sub-goals.
Closing a goal (replacing it with a completed sub-proof)
corresponds to $r(g) = \emptyset$.
\end{definition}

\begin{definition}[Construction functor]\label{def:confunctor}
The \emph{construction functor}
$\mathcal{C} \colon \mathbf{PP} \to \mathbf{PP}$
is defined by: for each partial proof~$P$, $\mathcal{C}(P)$
is the partial proof obtained by applying the participant's
suggested refinement steps to the open goals of~$P$.
\end{definition}

\begin{remark}
The construction functor is not guaranteed to produce a
complete proof.  It may reduce the number of open goals,
leave them unchanged, or even increase them (if a step
introduces new sub-goals).  The key property is that
trust is preserved.
\end{remark}

\subsection{Witnesses and construction plans}

\begin{definition}[Witness]\label{def:witness}
A \emph{witness} for a goal $(c, \exists x.\, \varphi(x))$
is a term~$t$ such that $\varphi(t)$ is valid.  The
participant suggests witnesses; the solver validates them.
\end{definition}

\begin{definition}[Construction plan]\label{def:conplan}
A \emph{construction plan} is a sequence of refinement steps
$r_1, r_2, \ldots, r_k$ together with witnesses for each
existential goal.  The plan is \emph{valid} if applying all
steps produces a complete proof.
\end{definition}

\begin{lemma}[Plan validity is decidable]\label{lem:planvalid}
Given a construction plan and the underlying proof obligations,
checking whether the plan produces a complete proof is decidable
(assuming the underlying logic is decidable for the given
fragment).
\end{lemma}
\begin{proof}
Each refinement step replaces a goal with sub-goals.  Checking
whether the resulting partial proof has no open goals is a
finite tree traversal.  If the logic is decidable, each
sub-goal can be checked.  Therefore plan validity is decidable.
\end{proof}

% ============================================================
%  Section 4 — The Construction Participant
% ============================================================
\section{The Construction Participant}\label{sec:participant}

\subsection{Architecture}

The $\CCP$ has three phases:
\begin{enumerate}
\item \emph{Plan generation.}  Given the current partial proof,
      the participant queries the Copilot channel for a
      construction plan.
\item \emph{Step suggestion.}  For each open goal, the
      participant suggests a refinement step (tactic).
\item \emph{Witness proposal.}  For existential goals, the
      participant proposes a witness term.
\end{enumerate}

\begin{definition}[Construction participant]\label{def:ccp}
The $\CCP$ is a triple $(\mathcal{P}, \tau, j)$ where:
\begin{itemize}
\item $\mathcal{P}$ is the participation function
      $\mathcal{P} \colon \mathbf{PP} \to \mathsf{Plan}$,
\item $\tau = \Tcopilot$ is the trust ceiling,
\item $j$ is the jurisdiction (all proposition kinds).
\end{itemize}
\end{definition}

\subsection{Plan generation}

The participant queries the Copilot channel with a prompt
encoding the current partial proof.  The prompt includes:
\begin{itemize}
\item The open goals (as propositions).
\item The names and types of available lemmas.
\item The proof history (steps taken so far).
\item A request for a ranked list of construction plans.
\end{itemize}

\subsection{Step suggestion}

For each open goal, the participant suggests a refinement
step.  The suggestion is a pair $(\mathsf{tactic}, \mathsf{args})$
where the tactic is one of: \texttt{intro}, \texttt{apply},
\texttt{split}, \texttt{induction}, \texttt{contradiction},
\texttt{exact}, \texttt{unfold}, \texttt{rewrite}.

\subsection{Witness proposal}

For existential goals, the participant proposes a witness.
The witness is a term in the language of the proof system.
The solver channel validates the witness by checking that
the instantiated formula is valid.

\begin{lemma}[Witness trust ceiling]\label{lem:witnessceil}
Every proposed witness has trust $\tleq \Tcopilot$.
\end{lemma}
\begin{proof}
Witnesses are derived from Copilot channel queries, which
have trust ceiling $\Tcopilot$.  The participant does not
elevate trust.
\end{proof}

% ============================================================
%  Section 5 — Implementation
% ============================================================
\section{Implementation}\label{sec:code}

\subsection{Construction plan generation}

\begin{lstlisting}[style=jugeo-python,
  caption={Construction plan generation via CopilotChannel},
  label={lst:conplan}]
from jugeo.evidence.channels import (
    CopilotChannel,
    ChannelConfiguration,
    EvidenceChannel,
    EvidenceRequest,
    EvidenceResponse,
    ChannelJurisdiction,
    ChannelMonitor,
)

class CopilotConstructionParticipant:
    """Copilot-guided participant in proof construction."""

    def __init__(self, channel: CopilotChannel):
        self.channel = channel
        self.trust_ceiling = channel.TRUST_CEILING
        self.jurisdiction = ChannelJurisdiction.for_channel(
            EvidenceChannel.COPILOT
        )

    def generate_plan(self, partial_proof, available_lemmas):
        prompt = self._build_plan_prompt(partial_proof, available_lemmas)
        raw = self.channel.query_llm(prompt, timeout_ms=12000)
        parsed = self.channel.parse_response()
        clamped = self.channel.apply_trust_ceiling()
        plans = self._extract_plans(clamped)
        return plans

    def suggest_step(self, goal, context):
        prompt = self._build_step_prompt(goal, context)
        raw = self.channel.query_llm(prompt, timeout_ms=8000)
        parsed = self.channel.parse_response()
        clamped = self.channel.apply_trust_ceiling()
        return self._parse_tactic(clamped)

    def propose_witness(self, existential_goal, context):
        prompt = self._build_witness_prompt(existential_goal, context)
        raw = self.channel.query_llm(prompt, timeout_ms=8000)
        parsed = self.channel.parse_response()
        clamped = self.channel.apply_trust_ceiling()
        return self._extract_witness(clamped)

    def _build_plan_prompt(self, proof, lemmas):
        goals = proof.get("open_goals", [])
        return (
            f"Proof has {len(goals)} open goals. "
            f"Available lemmas: {lemmas[:10]}. "
            f"Suggest a construction plan."
        )

    def _build_step_prompt(self, goal, ctx):
        return f"Goal: {goal}. Context: {ctx}. Suggest a tactic."

    def _build_witness_prompt(self, goal, ctx):
        return f"Existential goal: {goal}. Suggest a witness term."

    def _extract_plans(self, response):
        if isinstance(response, list):
            return response
        return [{"steps": [], "confidence": 0.5}]

    def _parse_tactic(self, response):
        return {"tactic": "apply", "args": response}

    def _extract_witness(self, response):
        return {"witness": response, "trust": self.trust_ceiling}
\end{lstlisting}

\subsection{Construction via CopilotAPIBridge}

\begin{lstlisting}[style=jugeo-python,
  caption={Proof construction with CopilotAPIBridge and JuGeoAPI},
  label={lst:conbridge}]
from jugeo.interfaces.api import (
    CopilotAPIBridge,
    APIEventLog,
    APIRequest,
    APIResponse,
    OperationKind,
    JuGeoAPI,
    APISession,
    RequestStatus,
)
from jugeo.evidence.channels import (
    CopilotChannel,
    ChannelConfiguration,
    EvidenceChannel,
)

class ConstructionOrchestrator:
    """Orchestrate proof construction with Copilot assistance."""

    def __init__(self, bridge: CopilotAPIBridge, api: JuGeoAPI):
        self.bridge = bridge
        self.api = api
        self.event_log = APIEventLog()

    def construct_with_copilot(self, coordinate, proposition):
        session = self.api.open_session("construction")
        hint = self.bridge.query(
            f"Construct proof for {proposition} at {coordinate}"
        )
        self.bridge.enforce_trust_ceiling()
        self.event_log.record({
            "event": "construction_hint",
            "hint": str(hint),
        })

        request = APIRequest(
            request_id="construct-001",
            operation=OperationKind.CONSTRUCT,
            coordinate=coordinate,
            proposition=proposition,
            budget=50000,
            deadline=60000,
            metadata={"copilot_hint": str(hint)},
        )
        session.register_request(request)
        result = self.api.construct(
            coordinate=coordinate,
            proposition=proposition,
        )

        self.event_log.record({
            "event": "construction_result",
            "success": result.is_successful() if result else False,
        })
        session.checkpoint()

        if result and result.has_residuals():
            residuals = self._handle_residuals(result, session)
            return {"result": result, "residuals": residuals}

        session.close()
        return {"result": result, "residuals": []}

    def _handle_residuals(self, result, session):
        residuals = []
        hint = self.bridge.query(
            f"Handle residuals from construction: {result}"
        )
        self.bridge.enforce_trust_ceiling()
        residuals.append({"hint": str(hint), "stats": self.bridge.stats()})
        return residuals
\end{lstlisting}

\subsection{Federated construction evidence}

\begin{lstlisting}[style=jugeo-python,
  caption={Federated evidence for construction validation},
  label={lst:confed}]
from jugeo.evidence.channels import (
    ChannelFederation,
    ChannelPool,
    ChannelRouter,
    ChannelMonitor,
    CopilotChannel,
    SolverChannel,
    RuntimeChannel,
    EvidenceChannel,
    EvidenceRequest,
    EvidenceResponse,
    EvidenceRecord,
    SupportRoute,
)

def federated_construction_validation(
    pool: ChannelPool,
    coordinate,
    proposition,
    witness=None,
):
    """Validate a construction step using federated evidence."""
    federation = ChannelFederation()
    router = ChannelRouter()
    monitor = ChannelMonitor()

    request = EvidenceRequest(
        request_id=f"con-val-{coordinate}",
        coordinate=coordinate,
        proposition=proposition,
        required_kind=None,
        proposition_kind="construction_validation",
        preferred_channel=EvidenceChannel.SOLVER,
        fallback_channels=[
            EvidenceChannel.COPILOT,
            EvidenceChannel.RUNTIME,
        ],
        deadline_ms=20000,
        budget=30000,
        metadata={"witness": str(witness)},
    )

    monitor.record_request(request)
    best = router.find_best_channel(request)
    all_capable = router.find_all_capable(request)

    evidence = federation.collect_all(request)
    merged = federation.merge_by_kind(evidence)
    federation.verify_no_escalation(merged)
    composed_trust = federation.compose_trust(merged)

    for ev in evidence:
        monitor.record_response(ev)

    return {
        "merged_evidence": merged,
        "composed_trust": composed_trust,
        "success_rate": monitor.success_rate(),
        "summary": monitor.summary(),
    }
\end{lstlisting}

\subsection{Step-by-step construction with health monitoring}

\begin{lstlisting}[style=jugeo-python,
  caption={Step-by-step construction with monitoring},
  label={lst:constep}]
from jugeo.kernel.health import (
    HealthStatus,
    HealthDimension,
    HealthIndicator,
    HealthReport,
)
from jugeo.evidence.channels import (
    ChannelPool,
    ChannelMonitor,
    CopilotChannel,
    EvidenceChannel,
    ChannelConfiguration,
)
from jugeo.interfaces.api import (
    JuGeoAPI,
    OperationKind,
    APIEventLog,
)

class StepByStepConstructor:
    """Build proofs step by step with health monitoring."""

    def __init__(self, pool: ChannelPool, api: JuGeoAPI):
        self.pool = pool
        self.api = api
        self.monitor = ChannelMonitor()

    def construct_steps(self, goals, max_steps=20):
        results = []
        for i, goal in enumerate(goals[:max_steps]):
            health = self._check_health()
            if health.overall_status == HealthStatus.UNHEALTHY:
                results.append({"goal": goal, "status": "skipped"})
                continue

            copilot = self.pool.get_channel("copilot")
            if copilot is None:
                results.append({"goal": goal, "status": "no_copilot"})
                continue

            raw = copilot.query_llm(
                f"Step {i}: prove {goal}", timeout_ms=10000
            )
            parsed = copilot.parse_response()
            clamped = copilot.apply_trust_ceiling()

            result = self.api.verify(
                coordinate=f"step-{i}",
                proposition=str(goal),
            )
            success = result.is_successful() if result else False
            results.append({
                "goal": goal,
                "suggestion": str(clamped),
                "verified": success,
            })
        return results

    def _check_health(self):
        copilot = self.pool.get_channel("copilot")
        indicators = []
        if copilot:
            alive = copilot.health_check()
            indicators.append(HealthIndicator(
                dimension=HealthDimension.COPILOT_CONNECTIVITY,
                status=(HealthStatus.HEALTHY if alive
                        else HealthStatus.UNHEALTHY),
                message="Copilot channel status",
                value=1.0 if alive else 0.0,
                timestamp=None,
                details={},
            ))
        overall = (HealthStatus.HEALTHY if all(
            i.is_healthy() for i in indicators
        ) else HealthStatus.UNHEALTHY)
        return HealthReport(
            overall_status=overall,
            indicators=indicators,
            generated_at=None,
            report_id="con-health",
            degradation_reasons=[],
            recovery_suggestions=[],
        )
\end{lstlisting}

% ============================================================
%  Section 6 — Formal Properties
% ============================================================
\section{Formal Properties}\label{sec:formal}

\begin{theorem}[Construction soundness]\label{thm:consound}
If the $\CCP$ suggests a construction plan~$\pi$ and every
step of~$\pi$ is validated by the solver channel, then the
completed proof is sound.
\end{theorem}
\begin{proof}
Each step of the plan is a refinement step $r_i$ that replaces
an open goal with sub-goals.  The solver validates each step
by checking that the sub-goals entail the parent goal.  Since
the solver is sound (at trust $\Tsolver$), each validated step
is correct.

The completed proof is the result of applying all validated steps.
By induction on the number of steps:
\begin{itemize}
\item \emph{Base case:} the empty plan leaves the proof unchanged.
\item \emph{Inductive step:} applying a validated step $r_i$ to
      a sound partial proof yields a sound partial proof with
      one fewer open goal (or additional sub-goals that are valid).
\end{itemize}
When all goals are closed, the proof is complete and sound.
\end{proof}

\begin{theorem}[Trust-ceiling preservation]\label{thm:trustceil}
Every suggestion produced by the $\CCP$ has trust $\tleq \Tcopilot$.
\end{theorem}
\begin{proof}
The participant queries the Copilot channel via \texttt{query\_llm}
and applies \texttt{apply\_trust\_ceiling}.  By the channel
contract, the response trust is at most $\Tcopilot$.  The
participant does not compose evidence from multiple channels
or elevate trust in any way.  Therefore every suggestion has
trust $\tleq \Tcopilot$.
\end{proof}

\begin{theorem}[Search-space reduction]\label{thm:searchred}
If the participant's plan reduces the number of open goals
from~$n$ to~$m < n$, then the remaining search space is
at most $b^m$ where $b$ is the branching factor, compared
to the original $b^n$.
\end{theorem}
\begin{proof}
The search space for a partial proof with $k$ open goals and
branching factor~$b$ is at most $b^k$ (each goal has at most
$b$ refinement options).  Reducing from~$n$ to~$m$ goals
reduces the search space from $b^n$ to $b^m$.  Since $m < n$,
$b^m < b^n$.
\end{proof}

\begin{theorem}[Witness completeness]\label{thm:witcomp}
For every decidable existential goal $(c, \exists x.\,\varphi(x))$,
if a witness exists, then the participant proposes at least
one candidate witness (in the idealised setting where the
language model is complete).
\end{theorem}
\begin{proof}
The participant queries the Copilot channel with the
existential goal.  In the idealised setting, the language
model enumerates all terms of the appropriate type and
checks $\varphi(t)$ for each.  Since a witness exists and
the goal is decidable, the model finds it.

In practice, the language model may miss witnesses.  The
theorem holds in the idealised setting; practical coverage
depends on the model's training data and inference budget.
\end{proof}

% ============================================================
%  Section 7 — Discussion
% ============================================================
\section{Discussion}\label{sec:discussion}

\paragraph{Practical witness completeness.}
Theorem~\ref{thm:witcomp} assumes an idealised language model.
In practice, the Copilot channel may fail to find a witness for
complex existential goals.  The construction participant mitigates
this by requesting multiple witnesses and ranking them by
confidence.

\paragraph{Cost of validation.}
Every participant suggestion requires solver validation, which
adds latency.  The trade-off is between construction speed
(with Copilot hints) and validation cost.  The participant
only suggests steps when the estimated speedup exceeds the
validation overhead.

\paragraph{Interaction with other advisors.}
The construction participant composes with the research advisor
(Paper~82) and the experiment advisor (Paper~82): the research
advisor suggests hypotheses that become construction goals, and
the experiment advisor designs validation experiments.

\paragraph{Residual handling.}
When construction produces residuals (open sub-goals), the
participant can be invoked recursively on each residual.
The recursion depth is bounded by the proof tree depth,
ensuring termination.

\paragraph{Trust and the validator.}
The participant operates at $\Tcopilot$, but the validated
proof has trust determined by the validator.  If the solver
validates all steps, the proof has trust $\Tsolver$.  The
participant's contribution is heuristic guidance, not trust.

% ============================================================
%  Section 8 — Related Work
% ============================================================
\section{Related Work}\label{sec:related}

\paragraph{LLM-guided proof construction.}
GPT-$f$~\cite{polu2020gptf} and AlphaProof~\cite{lample2022hypertree}
use language models to suggest proof steps.  Our participant
differs in operating within a trust-bounded channel architecture.

\paragraph{Proof assistants with AI.}
Lean Copilot~\cite{yang2019learning} and Isabelle
Sledgehammer~\cite{wenzel2014isabelle} integrate AI into
proof assistants.  Our approach is framework-independent and
mediated by the channel layer.

\paragraph{Witness generation.}
Property-based testing~\cite{claessen2000quickcheck} generates
witnesses by random search.  Our participant uses LLM-guided
search, which may find witnesses faster for structured goals.

\paragraph{Proof planning.}
Omega~\cite{basin1993rippling} and Clam~\cite{bundy1990clam}
use proof plans.  Our construction plans are generated by the
LLM rather than by hand-crafted heuristics.

\paragraph{Trust in AI proofs.}
The question of trusting AI-generated proofs is studied
by Ringer \etal{}~\cite{first2022diversity}.  Our trust lattice
provides a formal answer: AI suggestions are at $\Tcopilot$,
and trusted evidence requires solver validation.

% ============================================================
%  Section 9 — Conclusion
% ============================================================
\section{Conclusion}\label{sec:conclusion}

We have introduced the $\CCP$, a Copilot-guided agent that
participates in proof construction within the \JuGeo{}
framework.  The participant suggests construction plans,
refinement steps, and witnesses, all at $\Tcopilot$ trust.
Four theorems establish soundness (given solver validation),
trust-ceiling preservation, search-space reduction, and
witness completeness (in the idealised setting).

Future work includes adaptive plan generation based on
construction history, multi-step lookahead for witness
search, and integration with reinforcement learning for
tactic selection.

% ============================================================
%  Bibliography
% ============================================================
\begin{thebibliography}{25}

\bibitem{polu2020gptf}
S.~Polu and I.~Sutskever,
``Generative language modeling for automated theorem proving,''
\emph{arXiv preprint arXiv:2009.03393}, 2020.

\bibitem{lample2022hypertree}
G.~Lample, M.-A.~Lachaux, T.~Lavril, X.~Martinet,
A.~Hayat, G.~Ebner, A.~Rodriguez, and T.~Lacroix,
``HyperTree proof search for neural theorem proving,''
\emph{NeurIPS}, 2022.

\bibitem{yang2019learning}
K.~Yang and J.~Deng,
``Learning to prove theorems via interacting with proof
assistants,''
\emph{ICML}, 2019.

\bibitem{wenzel2014isabelle}
M.~Wenzel,
``Isabelle/jEdit --- a prover IDE within the PIDE framework,''
\emph{AISC/Calculemus/MKM}, 2012.

\bibitem{claessen2000quickcheck}
K.~Claessen and J.~Hughes,
``QuickCheck: a lightweight tool for random testing of
Haskell programs,''
\emph{ICFP}, 2000.

\bibitem{basin1993rippling}
D.~Basin and T.~Walsh,
``A calculus for and termination of rippling,''
\emph{Journal of Automated Reasoning}, vol.~16,
pp.~147--180, 1996.

\bibitem{bundy1990clam}
A.~Bundy, A.~Stevens, F.~van Harmelen, J.~Ireland,
A.~Smaill, and T.~Walsh,
``Rippling: a heuristic for guiding inductive proofs,''
\emph{Artificial Intelligence}, vol.~62, pp.~185--253, 1993.

\bibitem{first2022diversity}
E.~First, M.~Rabe, T.~Ringer, and Y.~Brun,
``Diversity-driven automated formal verification,''
\emph{ICSE}, 2022.

\bibitem{young2024jugeo}
H.~Young,
``Judgment geometry: a framework for trust-calibrated
verification,''
\emph{JuGeo Technical Reports}, Paper~1, 2024.

\bibitem{young2024channels}
H.~Young,
``Evidence channels in judgment geometry,''
\emph{JuGeo Technical Reports}, Paper~88, 2024.

\bibitem{young2024advisors}
H.~Young,
``Copilot advisors for heap, scope, import, and callable
analysis,''
\emph{JuGeo Technical Reports}, Papers~80--81, 2024.

\bibitem{young2024lifecycle}
H.~Young,
``Kernel lifecycle and health monitoring in judgment geometry,''
\emph{JuGeo Technical Reports}, Paper~93, 2024.

\bibitem{demoura2008z3}
L.~de~Moura and N.~Bj{\o}rner,
``Z3: an efficient SMT solver,''
\emph{TACAS}, 2008.

\bibitem{moura2021lean4}
L.~de~Moura and S.~Ullrich,
``The Lean~4 theorem prover and programming language,''
\emph{CADE}, 2021.

\bibitem{bansal2019holist}
K.~Bansal, S.~Loos, M.~Rabe, C.~Szegedy, and S.~Wilcox,
``HOList: an environment for machine learning of higher-order
logic theorem proving,''
\emph{ICML}, 2019.

\bibitem{rabe2021mathematical}
M.~N.~Rabe, D.~Lee, K.~Bansal, and C.~Szegedy,
``Mathematical reasoning via self-supervised skip-tree
training,''
\emph{ICLR}, 2021.

\bibitem{jiang2022thor}
A.~Q.~Jiang, W.~Li, S.~Tworber, and Y.~Wu,
``Thor: wielding hammers to integrate language models and
automated theorem provers,''
\emph{NeurIPS}, 2022.

\bibitem{li2021isarstep}
W.~Li, L.~Yu, Y.~Wu, and L.~C.~Paulson,
``IsarStep: a benchmark for high-level mathematical
reasoning,''
\emph{ICLR}, 2021.

\bibitem{clarke1999mc}
E.~M.~Clarke, O.~Grumberg, and D.~A.~Peled,
\emph{Model Checking}.
MIT Press, 1999.

\bibitem{ribeiro2020trust}
M.~T.~Ribeiro, S.~Singh, and C.~Guestrin,
``'Why should I trust you?': explaining the predictions
of any classifier,''
\emph{KDD}, 2016.

\end{thebibliography}

% ============================================================
%  Appendix
% ============================================================
\appendix

\section{Auxiliary Proofs}\label{app:proofs}

\begin{lemma}[Partial proof refinement is monotone]%
\label{lem:refmono}
If $P \to P'$ is a refinement step, then the number of
completed sub-proofs in $P'$ is at least the number in~$P$.
\end{lemma}
\begin{proof}
A refinement step closes one goal (adding a completed sub-proof)
or replaces one goal with sub-goals (the existing completed
sub-proofs are unchanged).  In either case, the count of
completed sub-proofs is non-decreasing.
\end{proof}

\begin{lemma}[Construction functor preserves identity]%
\label{lem:confuncid}
$\mathcal{C}(\mathrm{id}_P) = \mathrm{id}_{\mathcal{C}(P)}$
for every partial proof~$P$.
\end{lemma}
\begin{proof}
The identity refinement makes no changes.  Applying the
construction functor to an unchanged proof leaves it unchanged.
\end{proof}

\begin{lemma}[Construction functor preserves composition]%
\label{lem:confunccomp}
For refinement steps $r_1 \colon P \to P'$ and
$r_2 \colon P' \to P''$:
$\mathcal{C}(r_2 \circ r_1) = \mathcal{C}(r_2) \circ \mathcal{C}(r_1)$.
\end{lemma}
\begin{proof}
The construction functor applies the participant's suggestions
to each refinement.  Composing refinements and then applying
suggestions yields the same result as applying suggestions to
each refinement separately and composing, since the suggestions
depend only on the open goals at each step.
\end{proof}

\begin{lemma}[Witness validation is decidable]%
\label{lem:witdecid}
For a decidable existential goal and a proposed witness~$t$,
checking $\varphi(t)$ is decidable.
\end{lemma}
\begin{proof}
Substituting $t$ for~$x$ in $\varphi(x)$ yields a closed
formula $\varphi(t)$.  Since the goal is decidable, checking
$\varphi(t)$ is decidable.
\end{proof}

\begin{lemma}[Plan application terminates]%
\label{lem:planterm}
Applying a finite construction plan to a partial proof
terminates.
\end{lemma}
\begin{proof}
The plan has finitely many steps.  Each step is applied
once.  Therefore the application terminates in at most $k$
steps where $k$ is the plan length.
\end{proof}

\begin{lemma}[Validated plan is correct]%
\label{lem:planvalcorrect}
If every step of a plan is validated by the solver, then
the plan is correct.
\end{lemma}
\begin{proof}
Each validated step is correct (by solver soundness).
The composition of correct steps is correct (by induction).
Therefore the plan is correct.
\end{proof}

\begin{lemma}[Participant trust does not accumulate]%
\label{lem:trustnoaccum}
Applying $k$ participant suggestions in sequence does not
elevate trust above $\Tcopilot$.
\end{lemma}
\begin{proof}
Each suggestion has trust $\tleq \Tcopilot$.  Sequential
composition takes the minimum: $\min(\Tcopilot, \ldots,
\Tcopilot) = \Tcopilot$.
\end{proof}

\begin{lemma}[Residual goals are well-formed]%
\label{lem:residwf}
Every residual goal produced by a construction step is
syntactically well-formed.
\end{lemma}
\begin{proof}
Refinement steps are defined as functions from goals to
sets of goals.  The domain and range are the same set
(well-formed goals).  Therefore residuals are well-formed.
\end{proof}

\begin{lemma}[Federation trust for construction validation]%
\label{lem:fedtrust}
When federated construction validation (Listing~\ref{lst:confed})
collects evidence from channels $c_1, \ldots, c_n$, the composed
trust is $\inf_i \mathrm{trust}(c_i)$.
\end{lemma}
\begin{proof}
The federation's \texttt{compose\_trust} method computes the
infimum (meet) of the contributing trusts.  By the no-escalation
invariant, the composed trust cannot exceed any contributing
trust.  Therefore the result is $\inf_i \mathrm{trust}(c_i)$.
\end{proof}

\begin{proposition}[Full construction pipeline correctness]%
\label{prop:fullconcorrect}
The full construction pipeline (participant suggests plan,
solver validates steps, federation merges evidence) produces
a correct proof at trust $\Tsolver$ (assuming solver soundness).
\end{proposition}
\begin{proof}
The participant suggests at $\Tcopilot$.  The solver validates
each step at $\Tsolver$.  The validated proof inherits the
solver's trust.  By Theorem~\ref{thm:consound}, the proof is
sound.  By Lemma~\ref{lem:fedtrust}, federated validation
produces evidence at trust $\geq \Tcopilot$ (and at $\Tsolver$
when the solver participates).
\end{proof}

\begin{remark}[Extension to interactive construction]
The $\CCP$ is non-interactive: it proposes and the kernel
disposes.  An interactive variant would present suggestions
to a human user and await acceptance or rejection.  The
trust model extends naturally: human-accepted suggestions
would carry trust $\Tproof$ (human tier), while auto-accepted
suggestions remain at $\Tcopilot$ pending validation.
\end{remark}

\section{Extended Formal Development}\label{app:extended}

\subsection{Construction Algebra}

\begin{definition}[Construction algebra]\label{def:conalg}
A \emph{construction algebra} is a tuple
$(\mathcal{G}, \circ, e, \tau)$ where $\mathcal{G}$ is the
set of construction operations, $\circ$ is sequential
composition, $e$ is the identity operation (no-op), and
$\tau \colon \mathcal{G} \to \mathcal{T}$ assigns a trust
tier to each operation.
\end{definition}

\begin{lemma}[Construction algebra is a monoid]%
\label{lem:conalgmonoid}
$(\mathcal{G}, \circ, e)$ is a monoid.
\end{lemma}
\begin{proof}
\emph{Associativity:} sequential composition of construction
operations is function composition, which is associative.
\emph{Identity:} the no-op operation satisfies
$g \circ e = g$ and $e \circ g = g$ for all $g$.
\end{proof}

\begin{lemma}[Trust is a monoid homomorphism]%
\label{lem:trusthom}
The trust function $\tau$ satisfies
$\tau(g_1 \circ g_2) = \min(\tau(g_1), \tau(g_2))$.
\end{lemma}
\begin{proof}
The trust of a composed operation is the minimum of the
component trusts, since the composition cannot produce
evidence above the lowest component's ceiling.
\end{proof}

\subsection{Induction Principles for Construction}

\begin{definition}[Well-founded proof tree]\label{def:wfpt}
A proof tree is \emph{well-founded} if every descending
chain of goals is finite.  Equivalently, the tree has
finite depth.
\end{definition}

\begin{lemma}[Construction on well-founded trees terminates]%
\label{lem:confterm}
If the proof tree is well-founded and the participant's plan
has finite length, then the construction process terminates.
\end{lemma}
\begin{proof}
Each plan step either closes a goal (decreasing the number of
open goals) or replaces a goal with sub-goals at strictly
lower depth.  Since the tree is well-founded, the depth
decreases strictly, and the process terminates.
\end{proof}

\begin{lemma}[Structural induction on partial proofs]%
\label{lem:structind}
Any property~$P$ that holds for complete proofs (base case)
and is preserved by refinement steps (inductive case) holds
for all partial proofs reachable from a complete proof by
``un-refining'' (opening goals).
\end{lemma}
\begin{proof}
Standard structural induction on the tree structure of
partial proofs.  The base case is a leaf (complete sub-proof).
The inductive case extends to internal nodes (open goals
with sub-trees).
\end{proof}

\subsection{Channel Interaction Properties}

\begin{lemma}[Copilot query idempotence for caching]%
\label{lem:cacheidemp}
If the same prompt is sent to the Copilot channel twice
(with caching enabled), the second call returns the cached
response, which has the same trust as the first.
\end{lemma}
\begin{proof}
The cache key is the canonical form of the prompt.  If the
prompt matches, the cached response is returned without
querying the LLM.  The cached response has the same trust
tier as the original, since trust is stored alongside the
response.
\end{proof}

\begin{lemma}[Solver validation is independent of participant]%
\label{lem:solverindep}
The solver channel validates a construction step independently
of how the step was generated.  Whether the step came from
the participant, a human, or another advisor does not affect
the validation result.
\end{lemma}
\begin{proof}
The solver checks the mathematical content of the step
(does the sub-goal entail the parent goal?) without
inspecting the provenance.  Therefore the validation is
source-independent.
\end{proof}

\begin{lemma}[Federation order independence]%
\label{lem:fedorderindep}
The order in which channels contribute evidence to the
federation does not affect the merged result (assuming
deterministic merge-by-kind).
\end{lemma}
\begin{proof}
The merge-by-kind operation groups evidence by kind and
takes the meet of trusts within each group.  Both grouping
and meet are order-independent (set union is commutative,
and meet is commutative).
\end{proof}

\begin{lemma}[Health check does not modify proof state]%
\label{lem:healthnomod}
Running a health check during construction does not modify
the partial proof or any open goals.
\end{lemma}
\begin{proof}
Health checks read the channel pool's state and monitor
statistics.  They do not invoke any construction operations
or modify the proof tree.
\end{proof}

\subsection{Complexity Bounds}

\begin{lemma}[Plan generation cost]%
\label{lem:plancost}
The cost of generating a construction plan is
$O(T_{\mathrm{llm}} + n \cdot T_{\mathrm{parse}})$
where $T_{\mathrm{llm}}$ is the LLM query latency, $n$ is
the number of open goals, and $T_{\mathrm{parse}}$ is the
per-goal parse time.
\end{lemma}
\begin{proof}
The participant makes one LLM query ($T_{\mathrm{llm}}$) and
parses the response for each of the $n$ goals
($n \cdot T_{\mathrm{parse}}$).  No other operations dominate.
\end{proof}

\begin{lemma}[Validation cost per step]%
\label{lem:valcost}
The cost of validating one construction step is
$O(T_{\mathrm{solver}})$ where $T_{\mathrm{solver}}$ is the
solver's query latency for the sub-goal.
\end{lemma}
\begin{proof}
Validation submits one formula to the solver and waits for
the response.  The cost is dominated by the solver's
processing time.
\end{proof}

\begin{proposition}[Total construction cost]%
\label{prop:totalcost}
The total cost of constructing a proof with $k$ steps is
$O(T_{\mathrm{llm}} + k \cdot (T_{\mathrm{parse}} + T_{\mathrm{solver}}))$.
\end{proposition}
\begin{proof}
One LLM query for plan generation, plus $k$ parse-and-validate
cycles.  The federation overhead per step is absorbed into
$T_{\mathrm{solver}}$ since it runs in parallel with solver
queries.
\end{proof}

\begin{remark}[Amortised cost with caching]
With response caching, repeated construction attempts for
similar goals amortise the LLM query cost.  The amortised
cost per step approaches $O(T_{\mathrm{solver}})$ as the
cache hit rate increases.
\end{remark}

\begin{remark}[Parallel validation]
When multiple construction steps are independent, their
validations can be parallelised.  With $p$ parallel solver
instances, the validation time reduces to
$O(k/p \cdot T_{\mathrm{solver}})$.
\end{remark}

\begin{lemma}[Construction participant is side-effect free]%
\label{lem:sefree}
The $\CCP$ does not modify the proof state.  It only reads
the current state and produces suggestions.
\end{lemma}
\begin{proof}
All operations of the participant are reads (query the proof
state, query the Copilot channel) followed by pure computations
(parse, filter, rank).  No write operation is performed on
the proof state.  Only the kernel modifies the proof state,
after validating the participant's suggestions.
\end{proof}

\begin{lemma}[Multiple participants compose]%
\label{lem:multipart}
Multiple construction participants can operate on the same
proof state concurrently without interference.
\end{lemma}
\begin{proof}
Since participants are side-effect free
(Lemma~\ref{lem:sefree}), concurrent reads do not
interfere.  Each participant produces an independent
suggestion list.  The kernel merges these lists and
validates each suggestion independently.
\end{proof}

\end{document}
